\begin{ExerciseBlock}
\ExerciseHeading{04}{Wärmestrahlung eines glühenden Metallblocks}{quantum_optics_thermal_radiation}
Ein Metallblock mit Oberfläche $A=\qty{0.5}{\metre\squared}$ wird in einem Hochtemperaturofen erhitzt. Sein mit einem Spektrometer aufgenommenes Emissionsspektrum ist unten dargestellt. Nähern Sie die glatte Einhüllende durch das Spektrum eines schwarzen Körpers an.
\begin{center}
\CachedTikzFigure{content/resources/chapters/quantum_optics/tikz/thermal_spectrum.tex}
\end{center}
\begin{enumerate}
\item Lesen Sie die Wellenlänge $\lambda_\mathrm{max}$ des Intensitätsmaximums ab. Nennen Sie eine mögliche Ursache für die Einbrüche.
\item Berechnen Sie die Temperatur des Metallblocks.
\item Bestimmen Sie die abgestrahlte Leistung des Metallblocks.
\item Um welchen Faktor ändert sich die Leistung, wenn die absolute Temperatur um 20 Prozent steigt?
\end{enumerate}
\end{ExerciseBlock}
\begin{Solution}
\par\medskip\noindent\textbf{Musterlösung}\par\nopagebreak[4]
\begin{enumerate}
\item Die Einhüllende hat ihr Maximum bei ungefähr $\lambda_\mathrm{max}\approx\qty{1.05}{\micro\metre}$. Schmale Einbrüche können beispielsweise durch Absorption auf dem optischen Weg entstehen.
\item Mit der Wienschen Konstanten $b\approx\qty{2.898e-3}{\metre\kelvin}$ folgt $T=b/\lambda_\mathrm{max}\approx\qty{2.76e3}{\kelvin}$.
\item Das Stefan-Boltzmann-Gesetz liefert $P=\sigma AT^4\approx\qty{1.64}{\mega\watt}$.
\item Für $T'=1{,}2T$ ist $P'/P=(T'/T)^4=1{,}2^4\approx2{,}07$. Mit dem abgelesenen Maximum erhält man $P'\approx\qty{3.4}{\mega\watt}$.
\end{enumerate}
\end{Solution}
